\apendice{Plan de Proyecto Software}

\section{Introducción}

En este anexo se recoge la planificación general seguida durante el desarrollo del Trabajo Fin de Grado, así como una valoración inicial de su viabilidad. El proyecto se ha abordado como una evolución de una aplicación ya existente para el diseño de diagramas Entidad-Relación, por lo que la planificación no partió de un desarrollo completamente nuevo, sino de una fase previa de análisis del sistema heredado.

El trabajo se organizó de forma incremental. En primer lugar se revisó el estado inicial del editor, se preparó el entorno de desarrollo y se analizaron los principales módulos de la aplicación. A partir de esa base, se fueron realizando tareas de estabilización, ampliación funcional, validación y redacción de la memoria.

La planificación del proyecto no debe entenderse como una secuencia totalmente cerrada desde el inicio, sino como una organización progresiva del trabajo. Algunas decisiones se fueron ajustando a medida que se comprendía mejor el funcionamiento interno del editor y las implicaciones de incorporar nuevos elementos del modelo Entidad-Relación.

\section{Planificación temporal}

El desarrollo del proyecto se dividió en varias fases principales. La primera fase estuvo centrada en la preparación del entorno de trabajo y en la toma de contacto con el proyecto heredado. Durante esta etapa se revisó la estructura del repositorio, se configuró el entorno local de desarrollo y se comprobaron los comandos principales de ejecución, construcción y pruebas.

Después se llevó a cabo una fase de análisis del sistema existente. En ella se estudió el funcionamiento del editor, especialmente la relación entre la representación visual gestionada con mxGraph y el modelo interno utilizado por la aplicación. Esta fase fue necesaria para comprender cómo se almacenaban las entidades, atributos y relaciones, y cómo se utilizaba esa información en la validación y en la generación de SQL.

Una vez comprendido el funcionamiento general del sistema, se abordó una fase de estabilización del núcleo del editor. El objetivo de esta etapa fue reducir algunas dependencias implícitas y corregir comportamientos que podían dificultar futuras ampliaciones. Esta parte del trabajo permitió disponer de una base más segura antes de incorporar nuevos elementos del modelo Entidad-Relación.

A continuación, el desarrollo se centró en ampliar las funcionalidades soportadas por el editor. En esta fase se fueron incorporando nuevos elementos y comportamientos relacionados con el modelo Entidad-Relación, prestando atención a que cada ampliación quedara integrada no solo en la interfaz, sino también en el modelo interno, las validaciones, la transformación relacional y la generación de SQL.

Paralelamente al desarrollo, se realizaron pruebas automáticas y revisiones manuales del comportamiento de la aplicación. Estas comprobaciones sirvieron para revisar casos concretos y para detectar posibles incoherencias entre lo que el usuario diseña en el lienzo y lo que la aplicación procesa internamente.

Finalmente, se inició la fase de redacción de la memoria y de los anexos. Esta fase se ha realizado de forma progresiva, comenzando por los capítulos cuyo contenido estaba más estabilizado y dejando pendientes de cierre aquellas partes que dependen de las últimas ampliaciones del proyecto.

% TODO: añadir tabla resumen de fases/milestones cuando esté cerrado el desarrollo.

\begin{table}[h]
    \centering
    \begin{tabular}{p{0.25\textwidth} p{0.60\textwidth}}
        \toprule
        \textbf{Fase} & \textbf{Descripción} \\
        \midrule
        Preparación y análisis inicial & Configuración del entorno, revisión del repositorio y toma de contacto con el proyecto heredado. \\
        Análisis del sistema heredado & Estudio de la arquitectura del editor, del modelo interno y de los módulos de validación y generación de SQL. \\
        Estabilización del núcleo & Revisión de dependencias internas, persistencia, sincronización entre modelo y vista y configuración de relaciones. \\
        Ampliación del modelo E/R & Incorporación progresiva de nuevos elementos del modelo Entidad-Relación y adaptación de las validaciones y transformaciones. \\
        Pruebas y revisión manual & Comprobación de funcionalidades mediante pruebas automáticas y revisión manual de casos de uso relevantes. \\
        Redacción y cierre & Redacción de la memoria, revisión de anexos y ajustes finales del proyecto. \\
        \bottomrule
    \end{tabular}
    \caption{Resumen de las fases principales del desarrollo.}
    \label{tab:fases-desarrollo}
\end{table}

\section{Estudio de viabilidad}

El estudio de viabilidad permite valorar si el proyecto puede llevarse a cabo con los recursos disponibles y dentro del contexto académico en el que se desarrolla. En este caso, se trata de un proyecto software realizado como Trabajo Fin de Grado, por lo que la viabilidad se analiza principalmente desde dos puntos de vista: el económico y el legal.

\subsection{Viabilidad económica}

Desde el punto de vista económico, el proyecto no requiere una inversión elevada en infraestructura. La aplicación se desarrolla como una herramienta web y puede ejecutarse en un entorno local utilizando un equipo personal y herramientas de desarrollo habituales.

La mayor parte del coste del proyecto está asociado al tiempo de trabajo dedicado al análisis, implementación, pruebas y redacción de la documentación. También se han utilizado herramientas y servicios de apoyo al desarrollo, como el editor de código, el sistema de control de versiones, la plataforma de alojamiento del repositorio y el entorno de redacción de la memoria.

% TODO: completar la estimación económica con horas dedicadas, coste/hora aproximado
% y herramientas utilizadas cuando se cierre la planificación final.

\subsection{Viabilidad legal}

Desde el punto de vista legal, el proyecto se apoya en tecnologías y librerías de uso habitual en el desarrollo web. Al tratarse de una evolución de una aplicación existente, es importante tener en cuenta tanto la licencia del proyecto heredado como las licencias de las dependencias utilizadas.

Durante el cierre del proyecto deberá revisarse que las librerías empleadas sean compatibles con el uso académico del trabajo y que se respeten las condiciones de sus respectivas licencias. También será necesario citar adecuadamente las herramientas principales utilizadas en el desarrollo y en la redacción de la memoria.

% TODO: revisar licencias concretas del proyecto y dependencias principales antes de cerrar este apartado.


